% !PRINTED 45
\weberpage{45}

The values $x$ and $\sqrt{f(x)}$ determined by (4), (5) then also
become functions of these variables, and (5) can be
differentiated.

Let us denote by $\delta$ the differentiation with respect to the
variables occurring in the coefficients of $P$, $Q$, where $x$
is regarded as constant, and by $dx$ the corresponding differential of $x$,
when $x$ is determined by (4) or (5) as a function of these
coefficients; the differentiation of (4) then gives
\[ 2P\delta P - 2Q\delta Qf(x) + F'(x)\,dx = 0, \]
from which, with the use of (5):
\weq{6}{\dfrac{Q\,\delta P - P\,\delta Q}{F'(x)} = \dfrac{dx}{2\sqrt{f(x)}},}
and since the numerator on the left side is of degree $2m-2$,
and $F(x)$ of degree $2m$, formula (2) can be applied,
and it follows that
\weq{7}{\sum \dfrac{dx}{\sqrt{f(x)}} = 0,}
when the sum is extended over all the roots of equation
 (5).

We apply this theorem to the simplest case, namely
 $m=2$, and then obtain the following theorem:

\gesperrt{If}
\weq{8}{x_1, \sqrt{f(x_1)}; \quad x_2, \sqrt{f(x_2)}; \quad x_3, \sqrt{f(x_3)}; \quad x_4, \sqrt{f(x_4)}}
\gesperrt{are four pairs of values for which some function of the
form}
\weq{9}{a + bx + cx^2 + \sqrt{f(x)}}
\gesperrt{vanishes, then these pairs of values are not entirely
independent of one another, but there subsists between
them the differential equation}
\weq{10}{\dfrac{dx_1}{\sqrt{f(x_1)}} + \dfrac{dx_2}{\sqrt{f(x_2)}} + \dfrac{dx_3}{\sqrt{f(x_3)}} + \dfrac{dx_4}{\sqrt{f(x_4)}} = 0.}

The dependence of these four pairs of values, and hence an integration
 of the differential equation (10), can however also be expressed
 in algebraic form, by setting the function
(9) equal to zero for $x = x_1, x_2, x_3, x_4$ and then eliminating
$a, b, c$. In this way one obtains the determinant equation:
%CONT>>
